\begin{table}[H]
  \centering
  \caption{Choques reportados en 2016 por los hogares marcados que caen y que no caen
  en pobreza, modelo principal, pobreza monetaria, 2013$\to$2016.}
  \label{tab:errores_choques}
  \footnotesize
  \begin{tabular}{l cc c}
    \toprule
    \textbf{Medida en 2016} & \textbf{VP} & \textbf{FP} & \textbf{Dif. estandarizada} \\
    \midrule
    N.\textsuperscript{o} de choques & 3.16 & 2.27 & 0.21\textsuperscript{*} \\
    Tuvo al menos un choque & 0.81 & 0.72 & 0.20\textsuperscript{*} \\
    N.\textsuperscript{o} de tipos de choque & 1.40 & 1.17 & 0.19\textsuperscript{*} \\
    Choque de salud & 0.38 & 0.32 & 0.11\textsuperscript{*} \\
    Choque económico & 0.21 & 0.19 & 0.04 \\
    Choque patrimonial & 0.07 & 0.08 & $-$0.03 \\
    Choque agropecuario & 0.53 & 0.48 & 0.10 \\
    Choque familiar & 0.30 & 0.21 & 0.20\textsuperscript{*} \\
    Choque severo & 0.02 & 0.02 & 0.05 \\
    Respuesta que erosiona activos & 0.26 & 0.29 & $-$0.06 \\
    Respuesta que protege activos & 0.34 & 0.31 & 0.08 \\
    Intensificó el trabajo & 0.20 & 0.22 & $-$0.03 \\
    Retiró hijos del colegio & 0.00 & 0.00 & $-$0.07 \\
    Redujo alimentos & 0.23 & 0.23 & $-$0.01 \\
    Se endeudó con entidad formal & 0.08 & 0.10 & $-$0.08 \\
    Se endeudó informalmente & 0.16 & 0.14 & 0.06 \\
    No hizo ningún ajuste & 0.19 & 0.18 & 0.02 \\
    \bottomrule
  \end{tabular}
  \begin{minipage}{0.95\textwidth}
    \vspace{4pt}
    \footnotesize \textit{Nota:} Hogares marcados por el modelo principal con el umbral F$_2$. VP: son pobres en 2016 ($n=566$); FP: no son pobres en 2016 ($n=1{,}527$). Choques reportados en la encuesta de 2016, posteriores a la predicción y no usados por el modelo; las variables binarias se reportan como proporciones. Diferencia estandarizada y \textsuperscript{*} como en la Tabla~\ref{tab:errores_nucleo}. Sin ponderar.
    Fuente: cálculos propios.
  \end{minipage}
\end{table}
